%! Graphing Rational Functions — Unit 5, Lesson 5.5 — Guided Notes
\documentclass[10pt]{article}
\usepackage{algebra2-article}
\usepackage{algebra2-boxes}
\newcommand{\TallMath}[1]{$\displaystyle #1\rule[-1.4em]{0pt}{3.2em}$}
% Standard coordinate grid + axes: {xmin}{xmax}{ymin}{ymax}
\newcommand{\lgrid}[4]{%
  \draw[step=1, linegray!40, very thin] (#1,#3) grid (#2,#4);
  \draw[->, semithick, charcoal] (#1,0)--(#2,0) node[below right, font=\scriptsize]{$x$};
  \draw[->, semithick, charcoal] (0,#3)--(0,#4) node[left, font=\scriptsize]{$y$};}

\begin{document}

\pageheader{Unit 5, Lesson 5.5}{Guided Notes}
\namedateperiod

\vspace{0.08in}

\begin{objectivebox}
\textbf{By the end of this lesson, I will be able to\dots}
\begin{itemize}\setlength{\itemsep}{2pt}
  \item \textbf{factor} a rational function and sort every restriction into a \emph{hole} or a \emph{vertical asymptote} (A2.F.2h);
  \item find the \textbf{horizontal asymptote} by comparing degrees, and the \textbf{intercepts} by computing --- not by reading a picture (A2.F.2a/g);
  \item build a \textbf{sign chart} whose boundary points are the $x$-intercepts and the vertical asymptotes, and use it to say where the graph runs \emph{above} and \emph{below} the $x$-axis;
  \item \textbf{assemble and describe the whole graph} from its equation, and recognize when a function is really a transformed $\frac1x$ from Lesson 5.4 (A2.F.1c).
\end{itemize}
\end{objectivebox}

\vspace{0.05in}

\begin{vocabbox}
\small
Fill in each term as we build it together.
\par\vspace{2pt}   % \par is required: \termblanklong's \noindent is a no-op mid-paragraph
\termblanklong{Factored form (and why it comes first)}
\termblanklong{The cancel test: hole or wall}
\termblanklong{Degree comparison}
\termblanklong{Sign chart \& boundary points}
\end{vocabbox}

\begin{hookbox}
All three graphs below have a vertical asymptote at $x=-2$, a vertical asymptote at $x=3$, an
$x$-intercept at $(1,0)$, and branches that flatten toward $y=0$ at the far left and far right.
\textbf{Exactly one of them is the graph of} $f(x)=\dfrac{x-1}{(x+2)(x-3)}$.
\begin{center}
\begin{tikzpicture}[scale=0.28]
  \begin{scope}
    \lgrid{-6}{6}{-4}{4}
    \draw[navy, dashed, semithick] (-2,-4)--(-2,4);
    \draw[navy, dashed, semithick] (3,-4)--(3,4);
    \begin{scope} \clip (-6,-4) rectangle (6,4);
      \draw[very thick, forest, domain=-6:-2.16, samples=110, smooth] plot (\x,{((\x)-1)/(((\x)+2)*((\x)-3))});
      \draw[very thick, forest, domain=-1.85:2.9, samples=130, smooth] plot (\x,{((\x)-1)/(((\x)+2)*((\x)-3))});
      \draw[very thick, forest, domain=3.1:6, samples=110, smooth] plot (\x,{((\x)-1)/(((\x)+2)*((\x)-3))});
    \end{scope}
    \fill[charcoal] (1,0) circle (3.2pt);
    \node[charcoal, font=\small\bfseries] at (0,-5.6){A};
  \end{scope}
  \begin{scope}[xshift=13.5cm]
    \lgrid{-6}{6}{-4}{4}
    \draw[navy, dashed, semithick] (-2,-4)--(-2,4);
    \draw[navy, dashed, semithick] (3,-4)--(3,4);
    \begin{scope} \clip (-6,-4) rectangle (6,4);
      \draw[very thick, forest, domain=-6:-2.16, samples=110, smooth] plot (\x,{-((\x)-1)/(((\x)+2)*((\x)-3))});
      \draw[very thick, forest, domain=-1.85:2.9, samples=130, smooth] plot (\x,{-((\x)-1)/(((\x)+2)*((\x)-3))});
      \draw[very thick, forest, domain=3.1:6, samples=110, smooth] plot (\x,{-((\x)-1)/(((\x)+2)*((\x)-3))});
    \end{scope}
    \fill[charcoal] (1,0) circle (3.2pt);
    \node[charcoal, font=\small\bfseries] at (0,-5.6){B};
  \end{scope}
  \begin{scope}[xshift=27cm]
    \lgrid{-6}{6}{-4}{4}
    \draw[navy, dashed, semithick] (-2,-4)--(-2,4);
    \draw[navy, dashed, semithick] (3,-4)--(3,4);
    \begin{scope} \clip (-6,-4) rectangle (6,4);
      \draw[very thick, forest, domain=-6:-2.4, samples=110, smooth] plot (\x,{((\x)-1)/((((\x)+2)*((\x)+2))*((\x)-3))});
      \draw[very thick, forest, domain=-1.6:2.98, samples=130, smooth] plot (\x,{((\x)-1)/((((\x)+2)*((\x)+2))*((\x)-3))});
      \draw[very thick, forest, domain=3.02:6, samples=110, smooth] plot (\x,{((\x)-1)/((((\x)+2)*((\x)+2))*((\x)-3))});
    \end{scope}
    \fill[charcoal] (1,0) circle (3.2pt);
    \node[charcoal, font=\small\bfseries] at (0,-5.6){C};
  \end{scope}
\end{tikzpicture}
\end{center}
\vspace{-2pt}
\begin{itemize}\setlength{\itemsep}{2pt}
  \item Your first guess: \blank{1.0cm} \quad Why? \writeline
  \item Every feature you were given --- both asymptotes, the intercept, the flattening --- is true of
        \emph{all three}. So the checklist alone cannot decide. What else would you need to know?
        \writeline
\end{itemize}
\vspace{-2pt}
What separates them is which \emph{side} of the $x$-axis the graph sits on in each region. That is
exactly what a \textbf{sign chart} answers --- and it is the one tool today that turns a list of
features into a picture. We will settle A, B, or C at the end of the notes.
\end{hookbox}

\begin{notesbox}{1. The Build --- Five Steps, Always in This Order}
In Lesson 5.4 the asymptotes were printed on the face of the equation. Starting today they are not. The
denominator factors into \emph{two} pieces, the numerator may cancel one of them, and the graph has to
be \textbf{assembled} rather than read.

\vspace{4pt}
\begin{center}
\renewcommand{\arraystretch}{1.35}
\begin{tabularx}{\linewidth}{@{}c l X@{}}
\hline
\textbf{\#} & \textbf{Step} & \textbf{What you write down} \\
\hline
1 & \textbf{Factor} everything & numerator and denominator, completely \\
2 & \textbf{Restrict, then cancel} & restrictions from the \emph{original} denominator; each one becomes a \blank{1.6cm} or a \blank{2.6cm} \\
3 & \textbf{Compare degrees} & the horizontal asymptote: $y=\blank{0.8cm}$, $y=$ ratio, or \blank{1.2cm} \\
4 & \textbf{Compute intercepts} & $x$-int from the simplified \blank{1.8cm}; $y$-int is $f(\blank{0.5cm})$ \\
5 & \textbf{Sign chart} & boundary points, one test value each, above or below \\
\hline
\end{tabularx}
\end{center}

\vspace{3pt}
\begin{center}
\fbox{\parbox{0.92\linewidth}{\centering\vspace{2pt}
\textbf{Factor it, and every feature is already in there.} The sign chart is what turns the list into a
picture.\vspace{2pt}}}
\end{center}
\end{notesbox}

\begin{notesbox}{2. Step 2 Up Close --- Hole or Wall?}
This is Lesson 5.0's \emph{cancel test}, and Lesson 5.1's algebra is what runs it.

\vspace{3pt}
\begin{itemize}[leftmargin=1.5em, itemsep=3pt]
  \item Restrictions always come off the \textbf{original} denominator, \emph{before} anything is
        cancelled. Simplify first and a restriction disappears --- along with the feature it names.
  \item A factor that \textbf{cancels} leaves a \blank{1.6cm}: a single missing point. Find its
        height by substituting the excluded value into the \blank{2.6cm} expression.
  \item A factor that \textbf{survives} in the denominator leaves a \blank{3.2cm}, with equation
        $x=\blank{1.4cm}$.
\end{itemize}

\vspace{3pt}
\textbf{Quick sort.} Factor, then label each restriction.
\begin{center}
\renewcommand{\arraystretch}{1.6}
\begin{tabularx}{\linewidth}{@{}l X X X@{}}
\hline
\textbf{Function} & \textbf{Restrictions} & \textbf{Hole at\dots} & \textbf{Vertical asymptote(s)} \\
\hline
\TallMath{\frac{x-4}{(x-4)(x+1)}} & \blank{2.4cm} & \blank{2.0cm} & \blank{2.0cm} \\
\TallMath{\frac{x+5}{(x-2)(x+3)}} & \blank{2.4cm} & \blank{2.0cm} & \blank{2.0cm} \\
\TallMath{\frac{x^2-1}{x^2-x-2}}  & \blank{2.4cm} & \blank{2.0cm} & \blank{2.0cm} \\
\hline
\end{tabularx}
\end{center}
\vspace{-2pt}
{\small (The third one: $\dfrac{(x-1)(x+1)}{(x-2)(x+1)}$.)}
\end{notesbox}

\begin{notesbox}{3. Step 3 Up Close --- The Degree Comparison}
Nothing new here either; this is 5.0's rule, and it is the \emph{end behavior} statement (A2.F.2g)
written as a line.

\vspace{3pt}
\begin{center}
\renewcommand{\arraystretch}{1.4}
\begin{tabularx}{\linewidth}{@{}l l X@{}}
\hline
\textbf{Case} & \textbf{Horizontal asymptote} & \textbf{Example} \\
\hline
top degree $<$ bottom degree & $y=\blank{1.0cm}$ & $\dfrac{x-1}{x^2-x-6}$ \\
top degree $=$ bottom degree & $y=$ \blank{3.6cm} & $\dfrac{2x^2-8}{x^2-x-6}$, so $y=\blank{0.8cm}$ \\
top degree $>$ bottom degree & \blank{3.0cm} & $\dfrac{x^3}{x^2-1}$ \\
\hline
\end{tabularx}
\end{center}

\vspace{3pt}
Two cautions worth writing down.
\begin{itemize}[leftmargin=1.5em, itemsep=3pt]
  \item Use the degrees of the \emph{whole} top and the \emph{whole} bottom --- not of one factor. It
        does not matter whether you compare before or after cancelling, because cancelling a common
        factor lowers \blank{3.4cm} degree by the same amount.
  \item ``No horizontal asymptote'' is a complete answer. The graph still has end behavior --- the arms
        just run off to $\pm\infty$ instead of flattening.
\end{itemize}
\end{notesbox}

\begin{notesbox}{4. Worked Example A --- Two Walls, and a Sign Chart}
\[ f(x)=\frac{x-1}{x^2-x-6} \]

\textbf{Step 1 --- factor.} \quad $f(x)=\dfrac{x-1}{(\rule{1.6cm}{0.4pt})(\rule{1.6cm}{0.4pt})}$

\vspace{3pt}
\textbf{Step 2 --- restrict and cancel.} Restrictions: $x\neq\blank{0.9cm}$ and $x\neq\blank{0.9cm}$.
Does anything cancel? \blank{1.0cm} \; So there are \blank{0.7cm} holes, and the vertical asymptotes
are \blank{1.4cm} and \blank{1.4cm}.

\vspace{3pt}
\textbf{Step 3 --- degrees.} Top degree \blank{0.7cm}, bottom degree \blank{0.7cm}, so the horizontal
asymptote is \blank{1.4cm}.

\vspace{3pt}
\textbf{Step 4 --- intercepts.} $x$-intercept: the numerator is zero at $x=\blank{0.8cm}$, so the point
is \blank{1.4cm}. \; $y$-intercept: $f(0)=\dfrac{-1}{-6}=\blank{1.0cm}$, so the point is
\blank{1.4cm}.

\vspace{4pt}
\textbf{Step 5 --- sign chart.} The boundary points are the $x$-intercept and the two vertical
asymptotes: $-2$, $1$, and $3$. They cut the axis into \blank{0.7cm} intervals. Test one value in each.

\vspace{2pt}
\renewcommand{\arraystretch}{1.5}
\begin{tabularx}{\linewidth}{|C{2.6cm}|C{1.6cm}|C{4.0cm}|C{1.6cm}|X|}
\hline
\rowcolor{forestbg}
\textbf{Interval} & \textbf{Test $x$} & \textbf{$\dfrac{x-1}{(x+2)(x-3)}$} & \textbf{Sign} &
\textbf{Above / below} \\ \hline
$x<-2$    & $-3$ & & & \\ \hline
$-2<x<1$  & $0$  & & & \\ \hline
$1<x<3$   & $2$  & & & \\ \hline
$x>3$     & $4$  & & & \\ \hline
\end{tabularx}

\vspace{4pt}
\textbf{Assemble it.} Far to the left the graph is just \blank{1.6cm} the $x$-axis, hugging $y=0$;
it dives to $-\infty$ at the wall $x=-2$, comes back from $+\infty$ on the other side, falls through
$(1,0)$, drops to $-\infty$ at $x=3$, and returns from $+\infty$ to flatten onto $y=0$ on the right.

\vspace{4pt}
\begin{center}
\begin{tikzpicture}[scale=0.42]
  \lgrid{-6}{7}{-4}{4}
  \draw[navy, dashed, semithick] (-2,-4)--(-2,4);
  \draw[navy, dashed, semithick] (3,-4)--(3,4);
  \node[navy, font=\scriptsize, above left] at (-2,3.2) {$x=-2$};
  \node[navy, font=\scriptsize, above right] at (3,3.2) {$x=3$};
  \node[navy, font=\scriptsize, below right] at (5.4,0) {$y=0$};
  \begin{scope}
    \clip (-6,-4) rectangle (7,4);
    \draw[very thick, forest, domain=-6:-2.16, samples=120, smooth] plot (\x,{((\x)-1)/(((\x)+2)*((\x)-3))});
    \draw[very thick, forest, domain=-1.85:2.9, samples=140, smooth] plot (\x,{((\x)-1)/(((\x)+2)*((\x)-3))});
    \draw[very thick, forest, domain=3.1:7, samples=120, smooth] plot (\x,{((\x)-1)/(((\x)+2)*((\x)-3))});
  \end{scope}
  \fill[charcoal] (1,0) circle (3.2pt) node[below right, font=\scriptsize]{$(1,0)$};
\end{tikzpicture}
\end{center}

\vspace{2pt}
\textbf{Now the question 5.4 could not ask.} Does this graph ever \emph{cross} its horizontal asymptote
$y=0$? Set $f(x)=0$: that happens exactly where the \blank{2.4cm} is zero --- at $x=\blank{0.8cm}$.
So yes, and the crossing point is the \blank{2.4cm}. In 5.4 the numerator was a nonzero constant and a
crossing was impossible; here it is not.

\vspace{3pt}
Domain: \blank{6.0cm} \quad {\small(three intervals, joined with $\cup$ --- an interval may never span
a wall)}
\end{notesbox}

\begin{notesbox}{5. Worked Example B --- A Hole, and an Old Friend}
\[ g(x)=\frac{x^2-4}{x^2-x-6} \]

\textbf{Steps 1--2.} \quad $g(x)=\dfrac{(\rule{1.4cm}{0.4pt})(\rule{1.4cm}{0.4pt})}%
{(\rule{1.4cm}{0.4pt})(\rule{1.4cm}{0.4pt})}$ \quad Restrictions: $x\neq\blank{0.8cm}$,
$x\neq\blank{0.8cm}$.

\vspace{3pt}
The factor \blank{1.4cm} cancels, so $x=\blank{0.8cm}$ is a \blank{1.2cm}; the factor
\blank{1.4cm} survives, so $x=\blank{0.8cm}$ is a \blank{2.6cm}. Simplified:
$g(x)=\blank{2.2cm}$.

\vspace{3pt}
\textbf{The hole's coordinates.} Substitute the cancelled value into the \emph{simplified} expression:
$\dfrac{-2-2}{-2-3}=\blank{1.0cm}$, so the hole is at \blank{1.8cm}.

\vspace{3pt}
\textbf{Step 3.} Degrees \blank{0.7cm} and \blank{0.7cm} --- equal --- so the HA is the ratio of the
leading coefficients, $y=\blank{1.0cm}$.

\vspace{3pt}
\textbf{Step 4.} $x$-intercept from the \emph{simplified} numerator: \blank{1.4cm}. \;
$y$-intercept: $g(0)=\dfrac{-4}{-6}=\blank{1.0cm}$.

\vspace{4pt}
\begin{center}
\fbox{\parbox{0.92\linewidth}{\centering\vspace{2pt}
\textbf{A trap worth naming.} The original numerator is zero at $x=2$ \emph{and} at $x=-2$. But $-2$ is
not an $x$-intercept --- it is the hole. Only the zeros of the \emph{simplified} numerator are
intercepts.\vspace{2pt}}}
\end{center}

\vspace{4pt}
\textbf{Step 5 --- sign chart.} Boundary points: the $x$-intercept $2$ and the wall $3$. \emph{The hole
is not a boundary point} --- the sign does not change there, because the graph on both sides of it is
the same simplified function.
\begin{center}
\renewcommand{\arraystretch}{1.5}
\begin{tabularx}{0.94\linewidth}{|C{2.6cm}|C{1.6cm}|C{2.6cm}|X|}
\hline
\rowcolor{forestbg}
\textbf{Interval} & \textbf{Test $x$} & \textbf{$\dfrac{x-2}{x-3}$} & \textbf{Above / below} \\ \hline
$x<2$     & $0$   & & \\ \hline
$2<x<3$   & $2.5$ & & \\ \hline
$x>3$     & $4$   & & \\ \hline
\end{tabularx}
\end{center}

\vspace{3pt}
\begin{center}
\begin{minipage}{0.52\linewidth}\centering
\begin{tikzpicture}[scale=0.36]
  \lgrid{-6}{8}{-4}{6}
  \draw[navy, dashed, semithick] (3,-4)--(3,6);
  \draw[navy, dashed, semithick] (-6,1)--(8,1);
  \node[navy, font=\scriptsize, above left] at (3,5.2) {$x=3$};
  \node[navy, font=\scriptsize, below left] at (-1.4,1) {$y=1$};
  \begin{scope}
    \clip (-6,-4) rectangle (8,6);
    \draw[very thick, forest, domain=-6:2.8, samples=130, smooth] plot (\x,{1+1/((\x)-3)});
    \draw[very thick, forest, domain=3.2:8, samples=130, smooth] plot (\x,{1+1/((\x)-3)});
  \end{scope}
  \draw[forest, very thick, fill=white] (-2,0.8) circle (4.0pt);
  \node[charcoal, font=\scriptsize, below] at (-2.2,0.4) {hole};
  \fill[charcoal] (2,0) circle (3.2pt);
\end{tikzpicture}
\end{minipage}%
\begin{minipage}{0.46\linewidth}
\small
\textbf{And now look again.} Simplify all the way and rewrite the numerator using the denominator
(Lesson 5.4's disguise):
\[ \frac{x-2}{x-3}=\frac{(x-3)+\rule{0.7cm}{0.4pt}}{x-3}=\rule{2.4cm}{0.4pt} \]
So the whole graph is the rational \emph{parent}, shifted \blank{1.2cm} $3$ and \blank{1.2cm}
$1$ --- with one point punched out at \blank{1.6cm}.

\vspace{3pt}
Domain: \blank{4.2cm}
\end{minipage}
\end{center}
\end{notesbox}

\begin{notesbox}{6. Reading the Finished Picture --- and Settling the Hook}
Whatever the equation, the assembled graph has to obey all of this:
\begin{itemize}[leftmargin=1.5em, itemsep=3pt]
  \item It \textbf{never} touches a vertical asymptote --- that input is not in the \blank{1.8cm}.
  \item It \textbf{may} cross a horizontal asymptote, and if it does, the crossing is at an
        \blank{2.4cm}. (A transformed parent from 5.4 never can; its numerator is a nonzero
        constant.)
  \item Domain and range are written as \blank{1.4cm} of intervals, because an interval may never
        span a wall.
  \item Increasing/decreasing intervals are reported \emph{branch by branch}, for the same reason.
\end{itemize}

\vspace{4pt}
\textbf{Back to the Hook.} Your sign chart for $f(x)=\dfrac{x-1}{(x+2)(x-3)}$ said: below, above,
below, above.
\begin{itemize}[leftmargin=1.5em, itemsep=2pt]
  \item Graph \textbf{B} is wrong because \writeline
  \item Graph \textbf{C} is wrong because \writeline
  \item The correct graph is \blank{1.0cm}.
\end{itemize}
\end{notesbox}

\begin{practicebox}
Build $p(x)=\dfrac{x^2-x-6}{x^2-4}$ completely. Factor first.
\begin{enumerate}[leftmargin=1.7em, itemsep=5pt]
  \item Factored: $p(x)=\dfrac{(\rule{1.4cm}{0.4pt})(\rule{1.4cm}{0.4pt})}%
        {(\rule{1.4cm}{0.4pt})(\rule{1.4cm}{0.4pt})}$ \quad Restrictions: $x\neq\blank{0.8cm}$,
        $x\neq\blank{0.8cm}$
  \item Hole at \blank{2.0cm}; \quad vertical asymptote \blank{1.6cm}.
  \item Horizontal asymptote \blank{1.6cm}. \quad $x$-intercept \blank{1.6cm}. \quad
        $y$-intercept \blank{1.6cm}.
  \item Sign chart --- boundary points \blank{2.4cm}. Test $x=0$: sign \blank{1.0cm}; \;
        test $x=2.5$: sign \blank{1.0cm}; \; test $x=4$: sign \blank{1.0cm}.
  \item Domain: \blank{5.4cm}
\end{enumerate}
\end{practicebox}

\end{document}
